\documentclass[11pt]{article}

\usepackage[hyperref]{acl}
\usepackage{times}
\usepackage{latexsym}
\usepackage[T1]{fontenc}
\usepackage[utf8]{inputenc}
\usepackage{microtype}
\usepackage{inconsolata}
\usepackage{booktabs}
\usepackage{graphicx}
\usepackage{amsmath}
\usepackage{url}

\title{Domain-Specific Hallucination Detection in Large Language Models: Project Update 2}

\author{Varun Chundru \and Debasmita Biswas \\
  Department of Computer Science \\
  Purdue University Fort Wayne \\
  \texttt{\{vchundru, biswd01\}@pfw.edu}}

\begin{document}
\maketitle

\begin{abstract}
Large Language Models (LLMs) frequently produce hallucinations---fluent but factually unsupported outputs---posing risks in high-stakes domains. We develop a hallucination detection framework combining Natural Language Inference (NLI) with domain-specific fine-tuning, targeting response-level binary classification on the HaluEval benchmark. In this update, we report results across five approaches: (1) a TF-IDF + Logistic Regression baseline achieving F1=0.21, (2) a retrieval similarity detector (F1=0.66), (3) a zero-shot NLI detector (F1=0.43), (4) a fine-tuned DeBERTa-v3-base model achieving F1=0.91, and (5) a Monte Carlo Dropout ensemble achieving F1=0.93. We further demonstrate that uncertainty-based selective prediction can reach near-perfect accuracy, that temperature scaling reduces calibration error by 36.8\%, and that cross-dataset transfer to SciFact reveals significant domain gap. Code and model weights are available at: \url{https://github.com/varunteja99/hallucination-detection-nlp}.
\end{abstract}

% ============================================================================
\section{Introduction}
% ============================================================================

Large Language Models have demonstrated remarkable generative capabilities, yet they frequently produce hallucinations: outputs that are fluent and seemingly plausible but factually incorrect or unsupported by the provided context \cite{ji2023survey}. This problem is particularly dangerous in high-stakes domains such as healthcare, legal analysis, and scientific research, where users may trust model outputs without independent verification.

Hallucinations manifest in two primary forms. \textit{Intrinsic hallucinations} directly contradict the source material, while \textit{extrinsic hallucinations} introduce claims that cannot be verified from the given context. Both types undermine the reliability of LLM-based systems.

This project frames hallucination detection as a binary classification problem: given a knowledge context and an LLM-generated response, predict whether the response is \textit{faithful} or \textit{hallucinated}. We approach this through Natural Language Inference, hypothesizing that a response contradicting its source knowledge is hallucinated.

In this update, we present: (1) a simple ML baseline demonstrating that surface-level features are insufficient, (2) retrieval-augmented verification using dense embeddings, (3) Monte Carlo Dropout uncertainty quantification showing that incorrect predictions exhibit 5.75$\times$ higher uncertainty, (4) temperature scaling that reduces Expected Calibration Error from 0.078 to 0.050, and (5) cross-dataset evaluation on SciFact using PubMedBERT retrieval revealing significant domain transfer challenges.

% ============================================================================
\section{Related Work}
% ============================================================================

We organize related work around three threads that directly inform our approach: the hallucination detection benchmarks we build upon, the NLI-based detection paradigm we adopt, and the pre-trained model architecture we fine-tune.

\subsection{Hallucination Benchmarks}

Our work builds directly on the HaluEval benchmark \cite{li2023halueval}, which provides 30,000 samples across QA, dialogue, and summarization tasks with human-annotated hallucination labels. HaluEval's design---pairing a knowledge context with a generated response and a binary label---naturally maps to an NLI framing, which motivates our approach.

Other hallucination benchmarks inform our evaluation strategy. FEVER \cite{thorne2018fever} established the paradigm of claim verification against evidence documents. SciFact \cite{wadden2020fact} enables domain-specific evaluation on scientific claims, which we use for cross-dataset transfer evaluation in this update.

\subsection{NLI-Based Hallucination Detection}

Our detection approach is grounded in the insight that hallucination detection can be cast as an NLI problem: if a generated response is \textit{not entailed} by the source knowledge, it is likely hallucinated.

\citet{honovich2022true} demonstrated that NLI models can evaluate factual consistency by decomposing outputs into atomic claims and checking entailment against source documents. Their TRUE framework showed that off-the-shelf NLI models provide a strong zero-shot signal for hallucination detection, which we confirm with our zero-shot baseline. However, their work also revealed that zero-shot NLI suffers from low recall---a finding we replicate (recall=0.31)---motivating our fine-tuning approach.

SelfCheckGPT \cite{manakul2023selfcheckgpt} detects hallucinations by sampling multiple responses and measuring consistency. Our Monte Carlo Dropout experiment draws on similar intuition: multiple stochastic forward passes through the same model can reveal prediction uncertainty.

\subsection{DeBERTa and Pre-trained Language Models for NLI}

We select DeBERTa-v3-base \cite{he2021deberta,he2023debertav3} as our backbone because of its strong performance on NLI benchmarks and its architectural innovations. DeBERTa introduced \textit{disentangled attention}, which separately encodes content and position information:
\[
A_{i,j} = H_i H_j^\top + H_i P_{i|j}^\top + P_{i|j} H_j^\top
\]
This decomposition captures content-to-content, content-to-position, and position-to-content interactions separately, enabling nuanced reasoning about premise-hypothesis relationships. DeBERTa-v3 further improved pre-training efficiency through replaced token detection (ELECTRA-style).

% ============================================================================
\section{Methodology}
% ============================================================================

\subsection{Task Formulation}

Given a knowledge context $k$, a question or prompt $q$, and a generated response $r$, we predict a binary label $y \in \{0, 1\}$ indicating whether $r$ is factual ($y=0$) or hallucinated ($y=1$).

\subsection{Model 1: TF-IDF + Logistic Regression Baseline}

We extract TF-IDF features (50,000 features, unigrams + bigrams, sublinear TF) from the concatenated knowledge-question-answer text and train a Logistic Regression classifier (C=1.0, saga solver). This baseline tests whether surface-level lexical patterns suffice for hallucination detection.

\subsection{Model 2: Zero-Shot NLI Detection}

We use \texttt{MoritzLaurer/DeBERTa-v3-base-mnli-fever-anli} without HaluEval-specific training, mapping contradiction $\rightarrow$ hallucinated and entailment/neutral $\rightarrow$ factual.

\subsection{Model 3: Fine-Tuned DeBERTa-v3-base}

We initialize from \texttt{microsoft/deberta-v3-base} (184M parameters) and add a binary classification head. Training uses AdamW (lr=$2 \times 10^{-5}$), batch size 8, 3 epochs, max length 512, float32 precision, on TESLA T4x2 (Kaggle).

\subsection{Retrieval-Augmented Verification}

We compute cosine similarity between knowledge and response embeddings using the \texttt{all-MiniLM-L6-v2} sentence transformer. The hypothesis is that hallucinated responses exhibit lower similarity to source knowledge.

\subsection{Monte Carlo Dropout Uncertainty}

We run 20 forward passes through the fine-tuned DeBERTa model with dropout enabled at inference time. The mean probability across passes serves as the prediction; the standard deviation captures uncertainty. This implements the $\mathcal{L}_{\text{calibration}}$ component from our proposal.

\subsection{Temperature Scaling}

We learn a scalar temperature $T$ on the validation set by minimizing NLL: $\text{NLL}(\text{logits}/T, \text{labels})$. Calibrated probabilities $\text{softmax}(\text{logits}/T)$ are evaluated using Expected Calibration Error (ECE).

\subsection{Cross-Dataset Evaluation on SciFact}

We evaluate our HaluEval-trained model on SciFact \cite{wadden2020fact} (1,109 claims, 5,183 abstracts). We use PubMedBERT (\texttt{pritamdeka/S-PubMedBert-MS-MARCO}) for domain-specific retrieval, then classify retrieved evidence-claim pairs with our fine-tuned DeBERTa.

% ============================================================================
\section{Experimental Setup}
% ============================================================================

\subsection{Dataset}

\begin{table}[h]
\centering
\small
\begin{tabular}{lrrr}
\toprule
\textbf{Task} & \textbf{Train} & \textbf{Val} & \textbf{Test} \\
\midrule
QA & 7,000 & 1,500 & 1,500 \\
Dialogue & 7,000 & 1,500 & 1,500 \\
Summarization & 7,000 & 1,500 & 1,500 \\
\midrule
\textbf{Total} & \textbf{21,000} & \textbf{4,500} & \textbf{4,500} \\
\bottomrule
\end{tabular}
\caption{HaluEval dataset splits (70/15/15, seed=42).}
\label{tab:dataset}
\end{table}

The TF-IDF baseline was evaluated on the HaluEval QA subset (10,000 pairs expanded to 20,000 balanced samples).

\subsection{Evaluation Metrics}

We report Accuracy, Precision, Recall, F1-score, and AUROC. F1 is our primary metric. For calibration, we report ECE with 15 bins.

% ============================================================================
\section{Results}
% ============================================================================

\subsection{Main Results}

Table~\ref{tab:main_results} presents results across all approaches.

\begin{table}[h]
\centering
\small
\begin{tabular}{lccccc}
\toprule
\textbf{Model} & \textbf{Acc} & \textbf{P} & \textbf{R} & \textbf{F1} & \textbf{AUC} \\
\midrule
TF-IDF + LogReg & 0.21 & 0.21 & 0.22 & 0.21 & -- \\
Retrieval Sim. & 0.49 & 0.49 & 1.00 & 0.66 & 0.38 \\
Zero-Shot NLI & 0.60 & 0.73 & 0.31 & 0.43 & 0.65 \\
Fine-Tuned DeBERTa & 0.91 & 0.88 & 0.95 & 0.91 & 0.98 \\
MC Dropout (20) & \textbf{0.93} & \textbf{0.93} & 0.93 & \textbf{0.93} & \textbf{0.98} \\
Calibrated DeBERTa & 0.91 & 0.88 & \textbf{0.95} & 0.91 & 0.98 \\
\bottomrule
\end{tabular}
\caption{Overall results on HaluEval. TF-IDF evaluated on QA subset only.}
\label{tab:main_results}
\end{table}

The TF-IDF baseline achieves only 21\% accuracy---\textit{worse than random chance} on a balanced binary task. This occurs because the shared knowledge and question tokens dominate the TF-IDF representation, drowning out the discriminative signal from the answer. Additionally, ChatGPT-generated hallucinations have stylistic patterns that are anti-correlated with the true labels. This confirms that bag-of-words features fundamentally cannot capture the semantic reasoning required for hallucination detection, strongly justifying the use of transformer models.

Retrieval similarity alone achieves F1=0.66 but with a degenerate threshold (0.96), effectively classifying nearly all samples as hallucinated (recall=1.0, precision=0.49). The similarity distributions for factual and hallucinated samples overlap heavily.

MC Dropout averaging over 20 forward passes improves F1 from 0.91 to 0.93, improving both precision (0.88 $\rightarrow$ 0.93) and accuracy (0.91 $\rightarrow$ 0.93). The ensemble effect of averaging stochastic predictions reduces noise.

\subsection{Ablation: Performance by Task Type}

\begin{table}[h]
\centering
\small
\begin{tabular}{lcccc}
\toprule
\textbf{Task} & \textbf{Acc} & \textbf{P} & \textbf{R} & \textbf{F1} \\
\midrule
QA & 0.97 & 0.98 & 0.96 & 0.97 \\
Dialogue & 0.82 & 0.75 & 0.91 & 0.82 \\
Summarization & 0.96 & 0.93 & 0.98 & 0.96 \\
\bottomrule
\end{tabular}
\caption{Fine-tuned DeBERTa performance by task type.}
\label{tab:task_results}
\end{table}

QA achieves the highest F1 (0.97) due to clear factual grounding. Dialogue proves most challenging (F1=0.82) with lower precision (0.75), as conversational hallucinations involve more nuanced, implicit claims.

\subsection{Uncertainty Analysis}

MC Dropout uncertainty reveals a strong correlation with prediction correctness. Incorrect predictions exhibit a mean uncertainty of 0.145, which is \textbf{5.75$\times$ higher} than correct predictions (mean uncertainty 0.025). This confirms that the model reliably flags its own errors through uncertainty estimation.

The selective prediction curve shows that by filtering out uncertain predictions, accuracy approaches 100\%. At 5\% coverage (retaining only the most confident predictions), accuracy reaches 1.00; at 50\% coverage, accuracy remains above 0.99. This has practical implications: a human reviewer could focus exclusively on uncertain predictions, reducing review burden while maintaining high reliability.

\subsection{Calibration Results}

Temperature scaling reduced ECE from 0.078 to 0.050, a 36.8\% improvement. Before calibration, the model is overconfident---high-confidence predictions are correct less often than their confidence suggests. After calibration, predicted confidence more closely tracks actual accuracy. Temperature scaling does not change predictions or F1; it adjusts probability calibration for more reliable confidence estimates.

\subsection{Cross-Dataset Evaluation: SciFact}

\begin{table}[h]
\centering
\small
\begin{tabular}{lccccc}
\toprule
\textbf{Evaluation} & \textbf{Acc} & \textbf{P} & \textbf{R} & \textbf{F1} & \textbf{AUC} \\
\midrule
HaluEval (in-domain) & 0.91 & 0.88 & 0.95 & 0.91 & 0.98 \\
SciFact (cross-dataset) & 0.34 & 0.34 & 0.98 & 0.50 & 0.40 \\
\bottomrule
\end{tabular}
\caption{Cross-dataset transfer: HaluEval-trained model on SciFact (693 labeled claims).}
\label{tab:scifact}
\end{table}

Our HaluEval-trained model performs poorly on SciFact (F1=0.50), classifying 686 of 693 labeled claims as hallucinated regardless of true label (SUPPORT: 456, CONTRADICT: 237). This reveals a significant domain gap: the model has learned HaluEval-specific patterns rather than general factual consistency reasoning. Scientific claims use different vocabulary, reasoning patterns, and evidence structures. Domain-specific fine-tuning on scientific data would be required for effective scientific hallucination detection.

% ============================================================================
\section{Progress Since Update 1}
% ============================================================================

\begin{enumerate}
    \item \textbf{Addressed professor feedback:} Revised related work to connect directly to our methods. Added DeBERTa architecture description. Added simple ML baseline.
    \item \textbf{TF-IDF baseline:} 21\% accuracy on HaluEval QA confirms lexical features are anti-predictive.
    \item \textbf{Retrieval-augmented verification:} Dense similarity is insufficient as standalone detector (F1=0.66, degenerate threshold).
    \item \textbf{MC Dropout uncertainty:} Improved F1 from 0.91 to 0.93; incorrect predictions show 5.75$\times$ higher uncertainty.
    \item \textbf{Temperature scaling:} ECE reduced 36.8\% (0.078 $\rightarrow$ 0.050).
    \item \textbf{SciFact cross-dataset evaluation:} PubMedBERT retrieval + DeBERTa classification reveals significant domain gap (F1 drops 0.91 $\rightarrow$ 0.50).
    \item \textbf{Model deployment:} Fine-tuned model on HuggingFace (\url{varunchundru/hallucination-detector-deberta}).
\end{enumerate}

% ============================================================================
\section{Ideas That Didn't Work Out}
% ============================================================================

\paragraph{fp16 Mixed Precision Training.} Caused NaN values due to DeBERTa's disentangled attention. Resolved by using float32.

\paragraph{Retrieval Similarity as Standalone Detector.} Similarity distributions for factual and hallucinated samples overlap heavily, yielding a degenerate threshold that classifies nearly everything as hallucinated.

\paragraph{TF-IDF on Full Concatenated Text.} Achieves 21\%---below random chance---because shared context tokens dominate and ChatGPT-generated hallucinations have anti-correlated style patterns.

\paragraph{Cross-Dataset Transfer Without Fine-Tuning.} Model classifies 99\% of SciFact claims as hallucinated, indicating overfitting to HaluEval-specific patterns.

\paragraph{Span-Level Detection and Large Models.} Both deferred due to scope and compute constraints. Response-level classification with DeBERTa-base (184M) achieves strong results.

% ============================================================================
\section{Next Steps}
% ============================================================================

\begin{enumerate}
    \item \textbf{Context ablation:} Performance with/without knowledge context.
    \item \textbf{Training data scaling:} Learning curves at 10\%, 25\%, 50\%, 100\%.
    \item \textbf{Ensemble methods:} Combine DeBERTa + retrieval similarity + uncertainty.
    \item \textbf{Phase 3 (DPO):} If time permits, explore Direct Preference Optimization; otherwise Future Work.
    \item \textbf{Final presentation:} 20-minute presentation.
\end{enumerate}

\paragraph{Code Repository} \url{https://github.com/varunteja99/hallucination-detection-nlp}

\bibliography{custom}

\end{document}
